\section{AIME I 2016}
        \begin{enumerate}
        	\item (\textit{AIME I 2016}) For $-1<r<1$, let $S(r)$ denote the sum of the geometric series 
 \[12+12r+12r^2+12r^3+\cdots .\]  Let $a$ between $-1$ and $1$ satisfy $S(a)S(-a)=2016$. Find $S(a)+S(-a)$. 


 
 \textbf{Answer: }336



\textbf{Solution 1}

The sum of an infinite geometric series is $\frac{a}{1-r}\rightarrow \frac{12}{1\mp a}$. The product $S(a)S(-a)=\frac{144}{1-a^2}=2016$. $\frac{12}{1-a}+\frac{12}{1+a}=\frac{24}{1-a^2}$, so the answer is $\frac{2016}{6}=\boxed{336}$.

 



\textbf{Solution 2}

There is a clear infinite geometric series, so we use the geometric series formula $S = \frac{a}{1-r}$. However, adjusted for the context of this specific series, $S(a)=\frac{12}{1-a}$. 

 Plugging in the formula for both sum and product results in$\frac{12}{1-a} \times \frac{12}{1+a} = 2016$ for the product, and $\frac{12}{1-a} + \frac{12}{1+a}$ for the sum. This ensures both fractions have the same denominator, so that solving would be easier to work with.

 Upon first notice, dividing by $(1-a)(1+a)$ is equivalent to multiplying by 14. 

 Unsimping the top also results in a constant, as $12(1+a)+12(1-a)=12 + 12a + 12 - 12a = 24 + 0a = 24$.

 We can use the fact that dividing by $(1-a)(1+a)$ is equivalent to multiplying by 14 in this scenario, as we have a constant term, 24, in the numerator of the second fraction. Multiplying this by 14 results in $24 \times 14 = \boxed{336}$.

 


	\item (\textit{AIME I 2016}) Two dice appear to be normal dice with their faces numbered from $1$ to $6$, but each dice is weighted so that the probability of rolling the number $k$ is directly proportional to $k$. The probability of rolling a $7$ with this pair of dice is $\frac{m}{n}$, where $m$ and $n$ are relatively prime positive integers. Find $m+n$.


 
 \textbf{Answer: }071



\textbf{Solution}

It is easier to think of the dice as $21$ sided dice with $6$ sixes, $5$ fives, etc.  Then there are $21^2=441$ possible rolls. There are $2\cdot(1\cdot 6+2\cdot 5+3\cdot 4)=56$ rolls that will result in a seven.  The odds are therefore $\frac{56}{441}=\frac{8}{63}$.  The answer is $8+63=\boxed{071}$

 See also \href{https://artofproblemsolving.com/wiki/index.php?title=2006\_AMC\_12B\_Problems/Problem\_17}{2006 AMC 12B Problems/Problem 17}

 



\textbf{Solution 2}

Since the probability of rolling any number is 1, and the problem tells us the dice are unfair, we can assign probabilities to the individual faces. The probability of rolling $n$ is $\frac{n}{21}$ because $21=\frac{6 \cdot 7}{2}$Next, we notice that 7 can be rolled by getting individual results of 1 and 6, 2 and 5, or 3 and 4 on the separate dice. The probability that 7 is rolled is now $2(\frac{1}{21} \cdot \frac{6}{21}+\frac{2}{21} \cdot \frac{5}{21} + \frac{3}{21} \cdot \frac{4}{21})$ which is equal to $\frac{56}{441}=\frac{8}{63}$.  Therefore the answer is $8+63=\boxed{071}$~PEKKA

 



\textbf{Solution 3}

Since the probability of rolling a $1$ is $\frac{1}{21}$, the probability of rolling a $2$ is $\frac{2}{21}$ the probability of rolling a $3$ is $\frac{3}{21}$ and so on, we can make a chart of probabilities and add them together. Note that we only need the probabilities of $1$ and $6$, $2$ and $5$, and $3$ and $4$, and the rest is symmetry and the others are irrelevant. 

 We have: $2 \cdot (\frac{1}{21} \cdot \frac{2}{7}$ $+$ $\frac{2}{21} \cdot \frac{5}{21}$ $+$ $\frac{1}{7} \cdot \frac{4}{21})$ $=$ $2 \cdot (\frac{2}{147} + \frac{10}{441} + \frac{4}{147})$ = $2 \cdot \frac{4}{63} = \frac{8}{63}$. Therefore, the answer is $8 + 63$ = $\boxed{071}$



	\item (\textit{AIME I 2016}) A \textit{regular icosahedron} is a $20$-faced solid where each face is an equilateral triangle and five triangles meet at every vertex. The regular icosahedron shown below has one vertex at the top, one vertex at the bottom, an upper pentagon of five vertices all adjacent to the top vertex and all in the same horizontal plane, and a lower pentagon of five vertices all adjacent to the bottom vertex and all in another horizontal plane. Find the number of paths from the top vertex to the bottom vertex such that each part of a path goes downward or horizontally along an edge of the icosahedron, and no vertex is repeated.

                    \begin{figure}[H]
                        \centering
                        \includegraphics[width=0.5\textwidth]{Problem/AIME_I_2016_3_Problem_Figure_1.png}
                    \end{figure}
                    

 \textbf{Answer: }810



\textbf{Solution}

Think about each plane independently. From the top vertex, we can go down to any of 5 different points in the second plane. From that point, there are 9 ways to go to another point within the second plane: rotate as many as 4 points clockwise or as many 4 counterclockwise, or stay in place.  Then, there are 2 paths from that point down to the third plane.  Within the third plane, there are 9 paths as well (consider the logic from the second plane).  Finally, there is only one way down to the bottom vertex.  Therefore there are $5 \cdot 9 \cdot 2 \cdot 9 \cdot 1=\boxed{810}$ paths.

 



\textbf{Solution 2: Casework}

Assume an ant is on the top of this icosahedron. Note that the icosahedron has two pentagon planes and two points where the ant starts and ends. Also note that when the ant hits a vertex of the pentagon, there is only two ways to go down. When the ant ends up at the last vertex and is about to head down, there is only one way to go down.

 Case $1$: The ant move down to the first pentagon, then the next pentagon, and finally to the last point in a total of four moves; there are a total of $5 \cdot 2 \cdot 1 = 10$ ways to achieve this.

 Case $2$: The ant goes to the first pentagon and moves to another vertex in the same pentagon. The ant has a total of $8$ paths doing this since it may go through $1, 2, 3,$ or $4$ edges from either two directions when going across the edges of the pentagon (recall the ant may not repeat a move to another vertex) . Now the ant moves to the second pentagon below. The ant again has a total of $8$ moves if he wanders around the pentagon. Finally, the ant moves down after wandering to the last vertex. There is a total of $5 \cdot 8 \cdot 2 \cdot 8 = 640$ ways.

 Case $3$: Assume the ant goes to the first pentagon and wanders around. Then the ant goes the next pentagon and then heads directly to the last vertex. There are $5 \cdot 8 \cdot 2 = 80$ ways.

 Case $4$: Now let the ant go to the first pentagon and then go directly down to the next pentagon. The ant wanders around the second pentagon before heading to the last vertex. Like case $3$ above, there are $5 \cdot 2 \cdot 8 = 80$ ways.

 Adding up all four cases, we get $10 + 640 + 80 + 80 =\boxed{810}$ total paths, as desired.      ~skyscraper

 



\textbf{Solution 3}

Go to \href{https://artofproblemsolving.com/wiki/index.php/2020\_AMC\_10A\_Problems/Problem\_19}{2020 AMC 10A \#19}, and connect all of the centers of the faces on the dodecahedron to get the icosahedron. The answer is $\boxed{810}$.

 



\textbf{Solution 4}

We find the number of ways to go down from the top vertex to the pentagon with rotating around, then moving down to the second pentagon and rotating around a certain number of ways to be $640$ by saying there are $5$ ways to choose the first vertex to move down to, then $2$ choices for the direction to move in, and $4$ ways to chose how many vertices you will cross over, after multiplying by $2$ for the number of ways to get to the second pentagon and then by $2$ for direction and $4$ for number of vertices again we prove that the number of ways is $640$. 

 We then find the number of ways to rotate on the first pentagon and then go straight down to be $5$ times $2$ times $4$ times $2$, which is $80$.

 The number of ways to go straight from the top vertex to the second pentagon is $5$ times $2$ times $2$ times $4$, which is again $80$. 

 Lastly, we must find the number of ways to go straight from the top vertex to the bottom without rotating around any pentagon at all. This is done by multiplying $5$ by $2$, which is $10$. 

 We add these up to get our answer of $\boxed{810}$.~jamespudi

 


	\item (\textit{AIME I 2016}) A right prism with height $h$ has bases that are regular hexagons with sides of length $12$. A vertex $A$ of the prism and its three adjacent vertices are the vertices of a triangular pyramid. The dihedral angle (the angle between the two planes) formed by the face of the pyramid that lies in a base of the prism and the face of the pyramid that does not contain $A$ measures $60^\circ$. Find $h^2$.


 
 \textbf{Answer: }108



\textbf{Solution 1}

Let $B$ and $C$ be the vertices adjacent to $A$ on the same base as $A$, and let $D$ be the last vertex of the triangular pyramid.  Then $\angle CAB = 120^\circ$. Let $X$ be the foot of the altitude from $A$ to $\overline{BC}$. Then since $\triangle ABX$ is a $30-60-90$ triangle, $AX = 6$.  Since the dihedral angle between $\triangle ABC$ and $\triangle BCD$ is $60^\circ$, $\triangle AXD$ is a $30-60-90$ triangle and $AD = 6\sqrt{3} = h$. Thus $h^2 = \boxed{108}$.

 



\textbf{Solution 2}

Let $B$ and $C$ be the vertices adjacent to $A$ on the same base as $A$, and let $D$ be the last vertex of the triangular pyramid. Notice that we can already find some lengths. We have $AB=AC=12$ (given) and $BC=BD=\sqrt{144+h^2}$ by the Pythagorean Theorem. Let $M$ be the midpoint of $BC$. Then, we have $AM=6$ ($30-60-90$) triangles and $DM=\sqrt{36+h^2}$ by the Pythagorean Theorem. Applying the Law of Cosines, since $\angle AMD=60^{\circ}$, we get 
 \[h^2=36+h^2+36-\frac12 \cdot 12 \sqrt{36+h^2} \implies h^2=\boxed{108},\] as desired.

 -someone

 


	\item (\textit{AIME I 2016}) Anh read a book. On the first day she read $n$ pages in $t$ minutes, where $n$ and $t$ are positive integers. On the second day Anh read $n + 1$ pages in $t + 1$ minutes. Each day thereafter Anh read one more page than she read on the previous day, and it took her one more minute than on the previous day until she completely read the $374$ page book. It took her a total of $319$ minutes to read the book. Find $n + t$.


 
 \textbf{Answer: }053



\textbf{Solution 1}

Let $d$ be the number of days Anh reads for.  Because the difference between the number of pages and minutes Anh reads for each day stays constant and is an integer, $d$ must be a factor of the total difference, which is $374-319=55$.  Also note that the number of pages Anh reads is $dn+\frac{d(d-1)}{2}$.  Similarly, the number of minutes she reads for is $dt+\frac{d(d-1)}{2}$.  When $d$ is odd (which it must be), both of these numbers are multiples of $d$.  Therefore, $d$ must be a factor of $55$, $319$, and $374$.  The only such numbers are $1$ and $11$.  We know that Anh reads for at least $2$ days. Therefore, $d=11$.

 Using this, we find that she reads $55$ "additional" pages and $55$ "additional" minutes.  Therefore, $n=\frac{374-55}{11}=29$, while $t=\frac{319-55}{11}=24$.  The answer is therefore $29+24=\fbox{053}$.

 



\textbf{Solution 2}

We could see that both $374$ and $319$ are divisible by $11$ in the outset, and that $34$ and $29$, the quotients, are relatively prime. Both are the $average$ number of minutes across the $11$ days, so we need to subtract $\left \lfloor{\frac{11}{2}}\right \rfloor=5$ from each to get $(n,t)=(29,24)$ and $29+24=\boxed{053}$.

 



\textbf{Solution 3}

If we let $k$ be equal to the number of days it took to read the book, the sum of $n$ through $n+k$ is equal to $(2n+k)(k+1)=748$ Similarly, $(2t+k)(k+1)=638$ We know that both factors must be integers and we see that the only common multiple of $748$ and $638$ not equal to $1$ that will get us positive integer solutions for $n$ and $t$ is $11$. We set $k+1=11$ so $k=10$. We then solve for $n$ and $t$ in their respective equations, getting $2n+10=68$.  $n=29$  We also get $2t+10=58$.  $t=24$. Our final answer is $29+24=\boxed{053}$




\textbf{Solution 4}

Notice $374=34\cdot 11$ and $319=29\cdot 11$. Also, note the sum of an arithmetic series is $\frac{2n+k}{2} \cdot b$, where $n$ is our first term, $n+k$ is our final term, and $b$ is the number of terms. Since we know both sequences of $n$ and $t$ have the same length, and since $11$ is prime and shared by both $319$ and $374$, we deduce that $b=11$. Thus from here we know $2n+k=68$ and $2t+k=58$ by using our other factors $34$ and $29$. Finally, we add the two systems up and we get $2t+2k+2n=126$. But, notice that $k=b-1$, since the first term has $k=0$, and our last term has $k=b-1$. Plugging this back into our equation we get $2n+2t=106 \implies n+t=\boxed{053}$




\textbf{Solution 5}

We list two equations:\begin{align*}n+(n+1)+...+(n+k)&=374\\t+(t+1)+...+(t+k)&=319.\end{align*}
Subtracting the two, we get:
 \[(n-t)(k+1)=374-319=55.\]Manipulating the first and second equation, we get:\begin{align*}n(k+1)+\frac{k(k+1)}{2}&=374 \\t(k+1)+\frac{k(k+1)}{2}&=319.\end{align*}
We factor out the common factor $k+1$:\begin{align*}(k+1)\left(n+\frac{k}{2}\right)&=374 \\(k+1)\left(t+\frac{k}{2}\right)&=319.\end{align*}
Note that $374$ and $319$ have a GCD of $11,$ now combining this with our equation that $(n-t)(k+1)=55,$ we see that $k+1$ has to equal $11.$ Thus, we get:
 \[(n,t)=(29,24) \Rightarrow n+t=29+24=\boxed{053}.\]Note: We see that because n-t=5, it becomes impossible for n+k/2 and t+k/2 to both be multiples of 11. Thus, this satisfies our condition. Thus k+1 must be 11 to satisfy the common factor 11 constraint.\textbf{mathboy282}

 



\textbf{Solution 6 (Similar to 1, Diophantine Equations)}

Denote $d$ to be the number of days it takes Ana to read her book. Now, notice that the two sums

 
 \[\sum_{k=1}^{d} n+k-1 =374\]
 \[\text{and}\]
 \[\sum_{k=1}^{d} t+k-1 = 319\].

 Subtraction of the two sums yields 

 
 \[d(n-t) = 55\].

 The possible ordered pairing is $d = 5, n-t = 11$, $d = 11, n-t = 5$, $d = 55, n-t = 1$, or $d=1, n-t = 55$. Obviously the latter case doesn't work as obviously $373 +374 \neq 374$, and the second to last case won't work as $n$ would be negative. The first case won't work as $n$ is not an integer. The last case gives $(n,t) = (29,24)$ which satisfies the constraints so the answer is $\boxed{053}$



	\item (\textit{AIME I 2016}) In $\triangle ABC$ let $I$ be the center of the inscribed circle, and let the bisector of $\angle ACB$ intersect $\overline{AB}$ at $L$. The line through $C$ and $L$ intersects the circumscribed circle of $\triangle ABC$ at the two points $C$ and $D$. If $LI=2$ and $LD=3$, then $IC= \frac{p}{q}$, where $p$ and $q$ are relatively prime positive integers. Find $p+q$.


 
 \textbf{Answer: }013



\textbf{Solution 1}

Suppose we label the angles as shown below.
                    \begin{figure}[H]
                        \centering
                        \includegraphics[width=0.5\textwidth]{Solution/AIME_I_2016_6_Solution_Figure_1.png}
                    \end{figure}
                    As $\angle BCD$ and $\angle BAD$ intercept the same arc, we know that $\angle BAD=\gamma$. Similarly, $\angle ABD=\gamma$. Also, using $\triangle ICA$, we find $\angle CIA=180-\alpha-\gamma$. Therefore, $\angle AID=\alpha+\gamma=\angle DAI$, so $\triangle AID$ must be isosceles with $AD=ID=5$. Similarly, $BD=ID=5$. Then $\triangle DLB \sim \triangle ALC$, hence $\frac{AL}{AC} = \frac{3}{5}$. Also, $AI$ bisects $\angle LAC$, so by the Angle Bisector Theorem $\frac{CI}{IL} =\frac{AC}{AL}= \frac{5}{3}$. Thus $CI = \frac{10}{3}$, and the answer is $\boxed{013}$.

 



\textbf{Solution 2 (Incenter/Excenter)}

See the diagram in Solution 1.

 Let $\Delta ABC$'s $C$-excenter be $J_C$. Since $CI$ is an angle bisector, $\angle ACD = \angle BCD$, meaning that $D$ is the midpoint of arc $\overarc{AB}$. By the \href{https://artofproblemsolving.com/wiki/index.php?title=Incenter/Excenter\_Lemma}{Incenter/Excenter Lemma}, $DA=DI=DB=DJ_C=2+3=5$, and applying Power of a Point on circle $D$ gives $AL\cdot LB = LI\cdot LJ_C = 2(3+5) = 16$. Applying Power of a Point again on $\Delta ABC$'s circumcircle gives $AL\cdot LB = LC\cdot LD = 16$, and since $LD = 3$, $LC = \frac{16}{3}$. Thus $IC = LC-LI = \frac{16}{3}-2 = \frac{10}{3}$. We submit $10+3=\boxed{013}$.- NamelyOrange

 



\textbf{Solution 3}

WLOG assume $\triangle ABC$ is isosceles. Then, $L$ is the midpoint of $AB$, and $\angle CLB=\angle CLA=90^\circ$. Draw the perpendicular from $I$ to $CB$, and let it meet $CB$ at $E$. Since $IL=2$, $IE$ is also $2$ (they are both inradii). Set $BD$ as $x$. Then, triangles $BLD$ and $CEI$ are similar, and $\frac{2}{3}=\frac{CI}{x}$. Thus, $CI=\frac{2x}{3}$. $\triangle CBD \sim \triangle CEI$, so $\frac{IE}{DB}=\frac{CI}{CD}$. Thus $\frac{2}{x}=\frac{(2x/3)}{(2x/3+5)}$. Solving for $x$, we have:$x^2-2x-15=0$, or $x=5, -3$. $x$ is positive, so $x=5$. As a result, $CI=\frac{2x}{3}=\frac{10}{3}$ and the answer is $\boxed{013}$




\textbf{Solution 4}

WLOG assume $\triangle ABC$ is isosceles (with vertex $C$). Let $O$ be the center of the circumcircle, $R$ the circumradius, and $r$ the inradius. A simple sketch will reveal that $\triangle ABC$ must be obtuse (as an acute triangle will result in $LI$ being greater than $DL$) and that $O$ and $I$ are collinear. Next, if $OI=d$, $DO+OI=R+d$ and $R+d=DL+LI=5$. Euler gives us that $d^{2}=R(R-2r)$, and in this case, $r=LI=2$. Thus, $d=\sqrt{R^{2}-4R}$. Solving for $d$, we have $R+\sqrt{R^{2}-4R}=5$, then $R^{2}-4R=25-10R+R^{2}$, yielding $R=\frac{25}{6}$. Next, $R+d=5$ so $d=\frac{5}{6}$. Finally, $OC=OI+IC$ gives us $R=d+IC$, and $IC=\frac{25}{6}-\frac{5}{6}=\frac{10}{3}$. Our answer is then $\boxed{013}$.

 



\textbf{Solution 5}

Since $\angle{LAD} = \angle{BDC}$ and $\angle{DLA}=\angle{DCB}$, $\triangle{DLA}\sim\triangle{DBC}$. Also, $\angle{DAC}=\angle{BLC}$ and $\angle{ACD}=\angle{LCB}$ so $\triangle{DAC}\sim\triangle{BLC}$. Now we can call $AC$, $b$ and $BC$, $a$. By angle bisector theorem, $\frac{AD}{DB}=\frac{AC}{BC}$. So let $AD=bk$ and $DB=ak$ for some value of $k$. Now call $IC=x$. By the similar triangles we found earlier, $\frac{3}{ak}=\frac{bk}{x+2}$ and $\frac{b}{x+5}=\frac{x+2}{a}$. We can simplify this to $abk^2=3x+6$ and $ab=(x+5)(x+2)$. So we can plug the $ab$ into the first equation and get $(x+5)(x+2)k^2=3(x+2) \rightarrow k^2(x+5)=3$. We can now draw a line through $A$ and $I$ that intersects $BC$ at $E$. By mass points, we can assign a mass of $a$ to $A$, $b$ to $B$, and $a+b$ to $D$. We can also assign a mass of $(a+b)k$ to $C$ by angle bisector theorem. So the ratio of $\frac{DI}{IC}=\frac{(a+b)k}{a+b}=k=\frac{2}{x}$. So since $k=\frac{2}{x}$, we can plug this back into the original equation to get $\left(\frac{2}{x}\right)^2(x+5)=3$. This means that $\frac{3x^2}{4}-x-5=0$ which has roots -2 and $\frac{10}{3}$ which means our $CI=\frac{10}{3}$ and our answer is $\boxed{013}$.

 



\textbf{Solution 6}

Since $\angle BCD$ and $\angle BAD$ both intercept arc $BD$, it follows that $\angle BAD=\gamma$. Note that $\angle AID=\alpha+\gamma$ by the external angle theorem. It follows that $\angle DAI=\angle AID=\alpha+\gamma$, so we must have that $\triangle AID$ is isosceles, yielding $AD=ID=5$. Note that $\triangle DLA \sim \triangle DAC$, so $\frac{DA}{DL} = \frac{DC}{DA}$. This yields $DC = \frac{25}{3}$. It follows that $CI = DC - DI = \frac{10}{3}$, giving a final answer of $\boxed{013}$.

 



\textbf{Solution 7}

Let $I_C$ be the excenter opposite to $C$ in $ABC$. By the incenter-excenter lemma $DI=DC \therefore$ $LI_C=8,LI=2,II_C=10$. Its well known that $(I_C,I,L,C)=-1 \implies \dfrac{LI_C}{LI}.\dfrac{CI}{CI_C}=-1 \implies \dfrac{CI}{CI+10}=\dfrac{1}{4} \implies \boxed{CI=\dfrac{10}{3}}$.$\blacksquare$~Pluto1708

 Alternate solution: We can use the angle bisector theorem on $\triangle CBL$ and bisector $BI$ to get that $\frac{CI}{IL}=\frac{CI}{2}=\frac{BC}{BL}$. Since $\triangle CBL \sim \triangle ADL$, we get $\frac{BC}{BL}=\frac{AD}{DL}=\frac{5}{3}$. Thus, $CI=\frac{10}{3}$ and $p+q=\boxed{13}$.(\href{https://artofproblemsolving.com/community/c759169h1918283\_geometry\_problem}{https://artofproblemsolving.com/community/c759169h1918283\_geometry\_problem})

 



\textbf{Solution 8}

We can just say that quadrilateral $ADBC$ is a right kite with right angles at $A$ and $B$. Let us construct another similar right kite with the points of tangency on $AC$ and $BC$ called $E$ and $F$ respectively, point $I$, and point $C$. Note that we only have to look at one half of the circle since the diagram is symmetrical. Let us call $CI$ $x$ for simplicity's sake. Based on the fact that $\triangle BCD$ is similar to $\triangle FCI$ we can use triangle proportionality to say that $BD$ is $2\frac{x+5}{x}$. Using geometric mean theorem we can show that $BL$ must be $\sqrt{3x+6}$. With Pythagorean Theorem we can say that $3x+6+9=4{(\frac{x+5}{x})}^2$. Multiplying both sides by $x^2$ and moving everything to LHS will give you $3{x}^3+11{x}^2-40x-100=0$ Since $x$ must be in the form $\frac{p}{q}$ we can assume that $x$ is most likely a positive fraction in the form $\frac{p}{3}$ where $p$ is a factor of $100$. Testing the factors in synthetic division would lead $x = \frac{10}{3}$, giving us our desired answer $\boxed{013}$. ~Lopkiloinm

 



\textbf{Solution 9 (Cyclic Quadrilaterals)}


                    \begin{figure}[H]
                        \centering
                        \includegraphics[width=0.5\textwidth]{Solution/AIME_I_2016_6_Solution_Figure_1.png}
                    \end{figure}
                    Connect $D$ to $A$ and $D$ to $B$ to form quadrilateral $ACBD$. Since quadrilateral $ACBD$ is cyclic, we can apply Ptolemy's Theorem on the quadrilateral.

 Denote the length of $BD$ and $AD$ as $z$ (they must be congruent, as $\angle ABD$ and $\angle DAB$ are both inscribed in arcs that have the same degree measure due to the angle bisector intersecting the circumcircle at $D$), and the lengths of $BC$, $AC$, $AB$, and $CI$ as $a,b,c, x$, respectively.

 After applying Ptolemy's, one will get that:

 
 \[z(a+b)=c(x+5)\]

 Next, since $ACBD$ is cyclic, triangles $ALD$ and $CLB$ are similar, yielding the following equation once simplifications are made to the equation $\frac{AD}{CB}=\frac{AL}{BL}$, with the length of $BL$ written in terms of $a,b,c$ using the angle bisector theorem on triangle $ABC$:

 
 \[zc=3(a+b)\]

 Next, drawing in the bisector of $\angle BAC$ to the incenter $I$, and applying the angle bisector theorem, we have that:

 
 \[cx=2(a+b)\]

 Now, solving for $z$ in the second equation, and $x$ in the third equation and plugging them both back into the first equation, and making the substitution $w=\frac{a+b}{c}$, we get the quadratic equation:

 
 \[3w^2-2w-5=0\]

 Solving, we get $w=5/3$, which gives $z=5$ and $x=10/3$, when we rewrite the above equations in terms of $w$. Thus, our answer is $\boxed{013}$ and we're done.

 -mathislife52

 



\textbf{Solution 10(Visual)}


                    \begin{figure}[H]
                        \centering
                        \includegraphics[width=0.5\textwidth]{Solution/AIME_I_2016_6_Solution_Figure_3.png}
                    \end{figure}
                    




\textbf{Solution 11}

Let $AB=c,BC=a,CA=b$, and $x=\frac{a+b}{c}$. Then, notice that $\frac{CI}{IL}=\frac{a+b}{c}=x$, so $CI=IL\cdot{}x=2x$. Also, by the incenter-excenter lemma, $AD=BD=ID=IL+LD=5$. Therefore, by Ptolemy's Theorem on cyclic quadrilateral $ABCD$, $5a+5b=c(2x+5)$, so $5\left(\frac{a+b}{c}\right)=2x+5$, so $5x=2x+5$. Solving, we get that $x=\frac{5}{3}$, so $CI=\frac{10}{3}$ and the answer is $10+3=\boxed{013}$.

 



\textbf{Solution 12}

Perform a $\sqrt{bc}$ Inversion followed by a reflection along the angle bisector of $\angle BCA$.

 It's well known that
 \[AB \leftrightarrow \odot CBA \implies L \leftrightarrow D\]
 \[I \leftrightarrow I_A\]where $I_A$ is the $A-$excenter.

 Also by Fact 5, $DI_A = 5$.

 So,
 \[CL \cdot CD = CI \cdot CI_A\]
 \[\implies (CI + IL) \cdot (CI + ID) = (CI) \cdot (CI + II_A)\]
 \[\implies (CI + 2) \cdot (CI + 5) = (CI) \cdot (CI + 10)\]
 \[\implies 7CI +10= 10CI\]
 \[\implies CI = \boxed{\dfrac{10}{3}}.\blacksquare\]




\textbf{Solution 13}

Without loss of generality, let $\triangle ABC$ be isosceles. Note that by the incenter-excenter lemma, $DI = DA = DB.$ Hence, $DA=DB=5.$ Let the point of tangency of the incircle and $\overline{BC}$ be $F$ and the point of tangency of the incircle and $\overline{AC}$ be $E.$ We note that $\angle ALC = \angle BLC = 90^\circ$ and $LA=LB=4,$ which immediately gives $AE=BF=4.$ Applying the Pythagorean Theorem on $\triangle ALC$ and $\triangle IEC$ gives $2^2+x^2=y^2$ and $4^2+(2+y)^2 = (4+x)^2.$ Solving for $y$ gives us $y=\frac{10}{3}.$ Therefore, $IC = \frac{10}{3}$ so the answer is $\boxed{13}.$




\textbf{Solution 14 (Trig)}

Let $C_1\in AB$ be the point such that $IC_1\perp AB$, and let $D_1\in AB$ be defined similarly for $D$. We know that $\triangle IC_1L\sim\triangle DD_1L$, so by triangle ratios $DD_1=\frac{3}{2}r$, where $r$ is the inradius. Additionally, by cyclic quadrilaterals, we know that $\angle BAD=\angle DAB=\frac{\gamma}{2}$, where $\gamma$ is equivalent to $\angle ACB$. Thus $\triangle ADB$ is isosceles and $DD_1$ is the perpendicular bisector of the triangle, so $AD_1=\frac{c}{2}$. Since $\tan\left(\frac{\gamma}{2}\right)=\sqrt{\frac{(s-a)(s-b)}{s(s-c)}}$ from formulas (where $s$ is half the perimeter of $\triangle ABC$) and since $\tan\left(\frac{\gamma}{2}\right)=\frac{\frac{3}{2}r}{\frac{1}{2}c}$ from $\triangle ADB$, we can set up an equation:

 $\sqrt{\frac{(s-a)(s-b)}{s(s-c)}}=\frac{\frac{3}{2}r}{\frac{1}{2}c}\implies\frac{\sqrt{s(s-a)(s-b)(s-c)}}{s(s-c)}=\frac{3rs}{cs}\implies\frac{[ABC]}{s(s-c)}=\frac{3[ABC]}{cs}$
$\implies\frac{1}{s-c}=\frac{3}{c}\implies \frac{s}{c}=\frac{4}{3}$

 Let $C_2\in AB$ such that $CC_2\perp AB$. Then $CC_2=\frac{2[ABC]}{c}$. Using the area formula $[ABC]=rs$ and our fact from above yields $CC_2=\frac{8}{3}r$. We then notice that $\triangle CC_2L\sim\triangle IC_1L$, so if we let $x=CI$, by triangle ratios we find that $\frac{\frac{8}{3}r}{x+2}=\frac{r}{2}$, leading to $x=\frac{10}{3}$. Thus the answer is $10+3=\boxed{013}$.

 


	\item (\textit{AIME I 2016}) For integers $a$ and $b$ consider the complex number 
 \[\frac{\sqrt{ab+2016}}{ab+100}-\left(\frac{\sqrt{|a+b|}}{ab+100}\right)i.\] Find the number of ordered pairs of integers $(a,b)$ such that this complex number is a real number.


 
 \textbf{Answer: }103



\textbf{Solution}

We consider two cases:

 \textbf{Case 1:} $ab \ge -2016$.  

 In this case, if
 \[0 = \text{Im}\left({\frac{\sqrt{ab+2016}}{ab+100}-\left({\frac{\sqrt{|a+b|}}{ab+100}}\right)i}\right) = -\frac{\sqrt{|a+b|}}{ab+100}\]then $ab \ne -100$ and $|a + b| = 0 = a + b$.  Thus $ab = -a^2$ so $a^2 < 2016$.  Thus $a = -44,-43, ... , -1, 0, 1, ..., 43, 44$, yielding $89$ values. However since $ab = -a^2 \ne -100$,  we have $a \ne \pm 10$.  Thus there are $87$ allowed tuples $(a,b)$ in this case.

 \textbf{Case 2:} $ab < -2016$. 

 In this case, we want
 \[0 = \text{Im}\left({\frac{\sqrt{ab+2016}}{ab+100}-\left({\frac{\sqrt{|a+b|}}{ab+100}}\right)i}\right) = \frac{\sqrt{-ab-2016} - \sqrt{|a+b|}}{ab+100}\]Squaring, we have the equations $ab \ne -100$ (which always holds in this case) and 
 \[-(ab + 2016)= |a + b|.\]Then if $a > 0$ and $b < 0$, let $c = -b$.  If $c > a$, 
 \[ac - 2016 = c - a \Rightarrow (a - 1)(c + 1) = 2015 \Rightarrow (a,b) = (2, -2014), (6, -402), (14, -154), (32, -64).\]Note that $ab < -2016$ for every one of these solutions.  If $c < a$, then
 \[ac - 2016 = a - c \Rightarrow (a + 1)(c - 1) = 2015 \Rightarrow (a,b) = (2014, -2), (402, -6), (154, -14), (64, -32).\]Again, $ab < -2016$ for every one of the above solutions. This yields $8$ solutions.  Similarly, if $a < 0$ and $b > 0$, there are $8$ solutions.  Thus, there are a total of $16$ solutions in this case.

 Thus, the answer is $87 + 16 = \boxed{103}$.

 


	\item (\textit{AIME I 2016}) For a permutation $p = (a_1,a_2,\ldots,a_9)$ of the digits $1,2,\ldots,9$, let $s(p)$ denote the sum of the three $3$-digit numbers $a_1a_2a_3$ , $a_4a_5a_6$, and $a_7a_8a_9$. Let $m$ be the minimum value of $s(p)$ subject to the condition that the units digit of $s(p)$ is $0$. Let $n$ denote the number of permutations $p$ with $s(p) = m$. Find $|m - n|$.


 
 \textbf{Answer: }162



\textbf{Solution 1}

To minimize $s(p)$, the numbers $1$, $2$, and $3$ (which sum to $6$) must be in the hundreds places.  For the units digit of $s(p)$ to be $0$, the numbers in the ones places must have a sum of either $10$ or $20$. However, since the tens digit contributes more to the final sum $s(p)$ than the ones digit, and we are looking for the minimum value of $s(p)$, we take the sum's units digit to be $20$. We know that the sum of the numbers in the tens digits is $\sum\limits_{i=1}^9 (i) -6-20=45-6-20=19$. Therefore, $m=100\times6+10\times19+20=810$.

 To find $n$, realize that there are $3!=6$ ways of ordering the numbers in each of the places.  Additionally, there are three possibilities for the numbers in the ones place: $(4,7,9)$, $(5,7,8)$, and $(5,6,9)$.  Therefore there are $6^3\times3=648$ ways in total. $|m-n|=|810-648|=\fbox{162}$.

 



\textbf{Solution 2}

Like solution 1, to minimize \( s(p) \), we must have the hundreds digits to be 1, 2, and 3.

 Then, notice that the units digits must sum to 20, as it is impossible to get them to sum to 10 without 1, 2, or 3. Then, we choose the largest numbers to be in the units places to minimize \( s(p) \), resulting in us using 9, 7, and 4.

 Then, 5, 6, and 8 are in the tens slots, giving us three placeholder numbers, which are 159, 267, and 384, in which \( 159+267+384 = 810 \).

 To show 810 is the minimum possible notice that we must have the units digits to sum to 20, meaning we carry over a 2. Then, if we assume the middle is less than 21, say 20, then we arrive at a contradiction as we cannot make 20 using any of the others, and thus less than 20 as well is not possible, forcing us to carry over another 2. Then, if we carry over the 2, the smallest possible hundred digits are \( 1+2+3+2 = 8 \), and the smallest tens digit is 1 from 21. Thus, our smallest number must be 810.

 To find all \( s(p) = 810 \), we notice that the 100ds digits must be 1, 2, and 3 (from our proof earlier). In addition, we are to carry over a 2 from the tens places. (one can also see that if we have our 100ds digits \( a, b, \) and \( c \), then \( a+b+c = 6 \) where \( a \neq b \neq c \) and \( a,b,c > 0 \), meaning that we have only \( 1+2+3 \) that satisfies the condition). Then, say the middle digits are \( d, e, \) and \( f \). Then, since we are forced to carry over a 2 from the units digits, we have that \( d+e+f = 21-2 = 19 \). Then, we have \( d,e,f > 1 \) and \( d \neq e \neq f \neq 1,2,3 \). Then, we can do simple casework to obtain only three sets of numbers that work, those being, (4,6,9), (4,7,8), (5,6,8). We then double check the units digits to see if they truly do sum to 20, giving us (5,7,8), (5,6,9), and (4,7,9) respectively, which indeed sum to 20. Thus, we have three possible sums. Each of these three sums has \( 3! \times 3! \times 3! = 6^3 = 216 \) total permutations within digits, and thus there are \( 216 \times 3 = 648 \) total permutations (which is our value of \( n \)).

 Thus, \( |m-n| = |810 - 648| = \boxed{162} \).

 



\textbf{Solution 3(Similar to above solutions)}

We note that $a_3 + a_6 + a_9 \equiv 0\mod 10$. Since the maximum possible sum of these three digits if $9 + 8 + 7 = 24$ and the minimum is $1 + 2 + 3 = 6$ the only possible values for this sum is $10, 20$. Now, intuitively, to minimize $s(p)$ we need the hundred's digit of the final sum $s(p)$ to be minimized. Thus, the hundred's digits for each of the three summands must be minimized and we intuitively guess $1, 2, 3$. What I'm trying to say here is that we want the largest possible digits to be already used up as the unit's digits so that the smaller digits are left for the larger parts of the sums. So if we want to minimize $s(p)$ we want to maximize $a_3, a_6, a_9$ which in turn maximizes the final unit's digit sum. Out of $10, 20$ the maximum sum is $20$. We established earlier that we wanted to only use $1, 2, 3$ as the hundred's digits of each of the summands. Thus we can only use $4, 5, \cdots, 9$ as the only digits to sum to $20$. There are only $3$ ways namely $(6, 5, 9), (8, 5, 7), (9, 7, 4)$ for the unit's digits. There are $3$ ways to choose the unit's digits row just by choosing any of these $3$ ordered triplets and there are $3!$ ways to arrange each of the numbers within the triplets. Since we fixed $1, 2, 3$ as the hundred's digits there are $3!$ ways to arrange $1, 2, 3$ leaving only a fixed triplet for the ten's digits leaving only $3!$ ways to arrange these numbers. Ultimately, we have

 $3! \cdot 3! \cdot 3! \cdot 3 = 648$ ways to permute to $s(p)$. The hard part is done and now we just need to find the minimum sum which is very easy. We choose $(9, 7, 4)$ as the unit's digits and $(1, 2, 3)$ are our fixed hundred's digits leaving $(8, 6, 5)$ as the tens digits. So we have a sum of 

 $100(1 + 2 + 3) + 10(8 + 6 + 5) + (9 + 7 + 4) = 600 + 190 + 20 = 810$ as the minimum $s(p)$. Finally, we have

 $|810 - 648| = \boxed{162}$



	\item (\textit{AIME I 2016}) Triangle $ABC$ has $AB=40,AC=31,$ and $\sin{A}=\frac{1}{5}$. This triangle is inscribed in rectangle $AQRS$ with $B$ on $\overline{QR}$ and $C$ on $\overline{RS}$. Find the maximum possible area of $AQRS$.


 
 \textbf{Answer: }744



\textbf{Solution 1}

Note that if angle $BAC$ is obtuse, it would be impossible for the triangle to inscribed in a rectangle. This can easily be shown by drawing triangle ABC, where $A$ is obtuse. Therefore, angle A is acute. Let angle $CAS=n$ and angle $BAQ=m$. Then, $\overline{AS}=31\cos(n)$ and $\overline{AQ}=40\cos(m)$. Then the area of rectangle $AQRS$ is $1240\cos(m)\cos(n)$. By product-to-sum, $\cos(m)\cos(n)=\frac{1}{2}(\cos(m+n)+\cos(m-n))$. $\cos(m+n)=\sin(90-m-n)=\sin(BAC)=\frac{1}{5}$. The maximum possible value of $\cos(m-n)$ is 1, which occurs when $m=n$. Thus the maximum possible value of $\cos(m)\cos(n)$ is $\frac{1}{2}(\frac{1}{5}+1)=\frac{3}{5}$ so the maximum possible area of $AQRS$ is $1240\times{\frac{3}{5}}=\fbox{744}$.

 



\textbf{Solution 2}

We start by drawing a diagram;

 
                    \begin{figure}[H]
                        \centering
                        \includegraphics[width=0.5\textwidth]{Solution/AIME_I_2016_9_Solution_Figure_1.png}
                    \end{figure}
                    

 We know that $\sin \alpha = \frac{1}{5}$. Since $\sin \alpha = \cos (90- \alpha)$, 
 \[\cos (90- \alpha) = \frac{1}{5} \implies  \cos (\beta + \gamma) = \frac{1}{5}\]

 Using our angle sum identities, we expand this to $\cos \beta \cdot \cos \gamma - \sin \beta \cdot \sin \gamma = \frac{1}{5}$. We can now use the right triangle definition of cosine and sine to rewrite this equation as;

 
 \[\frac{l}{40} \cdot \frac{w}{31} - \frac{x}{40} \cdot \frac{y}{31} = \frac{1}{5} \implies lw- xy = 8 \cdot 31 \implies lw = xy + 31 \cdot 8\]

 Hang on; $lw$ is the area we want to maximize! Therefore, to maximize this area we must maximize $xy = 40 \sin \beta \cdot 31 \sin \gamma = 31 \cdot 40 \cdot \frac{1}{2}( \cos (\beta - \gamma) - \cos ( \beta + \gamma)) = 31 \cdot 20 \cdot (\cos(\beta-\gamma)-\frac{1}{5})$.Since $\cos(\beta-\gamma)$ is the only variable component of this expression, to maximize the expression we must maximize $\cos(\beta-\gamma)$. The cosine function has a maximum value of 1, so our equation evaluates to $xy = 31 \cdot 20 \cdot (1-\frac{1}{5}) = 31 \cdot 20 \cdot \frac{4}{5} = 31 \cdot 16$ (Note that at this max value, since $\beta$ and $\gamma$ are both acute, $\beta-\gamma=0 \implies \beta=\gamma$).

 Finally, $lw = xy + 31 \cdot 8 = 31 \cdot 16 + 31\cdot 8 = 31 \cdot 24 = \boxed{744}$




\textbf{Solution 3}

As above, we note that angle $A$ must be acute. Therefore, let $A$ be the origin, and suppose that $Q$ is on the positive $x$ axis and $S$ is on the positive $y$ axis. We approach this using complex numbers. Let $w=\text{cis} A$, and let $z$ be a complex number with $|z|=1$, $\text{Arg}(z)\ge 0^\circ$  and $\text{Arg}(zw)\le90^\circ$. Then we represent $B$ by $40z$ and $C$ by $31zw$. The coordinates of $Q$ and $S$ depend on the real part of $40z$ and the imaginary part of $31zw$. Thus
 \[[AQRS]=\Re(40z)\cdot \Im(31zw)=1240\left(\frac{z+\overline{z}}{2}\right)\left(\frac{zw-\overline{zw}}{2i}\right).\]We can expand this, using the fact that $z\overline{z}=|z|^2$, finding
 \[[AQRS]=620\left(\frac{z^2w-\overline{z^2w}+w-\overline{w}}{2i}\right)=620(\Im(z^2w)+\Im(w)).\]Now as $w=\text{cis}A$, we know that $\Im(w)=\frac15$. Also, $|z^2w|=1$, so the maximum possible imaginary part of $z^2w$ is $1$. This is clearly achievable under our conditions on $z$. Therefore, the maximum possible area of $AQRS$ is $620(1+\frac15)=\boxed{744}$.

 



\textbf{Solution 4 (With Calculus)}

Let $\theta$ be the angle $\angle BAQ$. The height of the rectangle then can be expressed as $h = 31 \sin (A+\theta)$, and the length of the rectangle can be expressed as $l = 40\cos \theta$. The area of the rectangle can then be written as a function of $\theta$, $[AQRS] = a(\theta) = 31\sin (A+\theta)\cdot 40 \cos \theta = 1240 \sin (A+\theta) \cos \theta$. For now, we will ignore the $1240$ and focus on the function $f(\theta) = \sin (A+\theta) \cos \theta = (\sin A \cos \theta + \cos A \sin \theta)(\cos \theta) = \sin A \cos^2 \theta + \cos A \sin \theta \cos \theta = \sin A \cos^2 \theta + \frac{1}{2} \cos A \sin 2\theta$.

 Taking the derivative, $f'(\theta) = \sin A \cdot -2\cos \theta \sin \theta + \cos A \cos 2\theta = \cos A \cos 2\theta - \sin A \sin 2\theta = \cos(2\theta + A)$. Setting this equal to $0$, we get $\cos(2 \theta + A) = 0 \Rightarrow 2\theta +A = 90, 270 ^\circ$. Since we know that $A+ \theta < 90$, the $270^\circ$ solution is extraneous. Thus, we get that $\theta = \frac{90 - A}{2} = 45 - \frac{A}{2}$.

 Plugging this value into the original area equation, $a(45 -  \frac{A}{2}) = 1240 \sin (45 - \frac{A}{2} + A) \cos (45 - \frac{A}{2}) = 1240\sin( 45+ \frac{A}{2})\cos(45 - \frac{A}{2})$. Using a product-to-sum formula, we get that: 
 \[1240\sin( 45+ \frac{A}{2})\cos(45 - \frac{A}{2}) =\]
 \[1240\cdot \frac{1}{2}\cdot(\sin((45 + \frac{A}{2}) + (45 -\frac{A}{2}))+\sin((45 +\frac{A}{2})-(45 - \frac{A}{2})))=\]
 \[620 (\sin 90^\circ + \sin A) = 620 \cdot \frac{6}{5} = \boxed{744}\].

 



\textbf{Solution 5}

Let $\alpha$ be the angle $\angle CAS$ and $\beta$ be the angle $\angle BAQ$. Then
 \[\alpha + \beta + \angle A = 90^\circ \Rightarrow \alpha + \beta = 90^\circ - \angle A\]
 \[\cos(\alpha + \beta) = \cos(90^\circ - \angle A)\]
 \[\cos(\alpha + \beta) = \sin(\angle A) = \frac{1}{5}\]
 \[\cos\alpha\cos\beta - \sin\alpha\sin\beta = \frac{1}{5}\]
 \[\cos\alpha\cos\beta - \sqrt{(1-\cos^2\alpha)(1-\cos^2\beta)} = \frac{1}{5}\]
 \[\cos\alpha\cos\beta - \sqrt{1-\cos^2\alpha-\cos^2\beta+\cos^2\alpha\cos^2\beta} = \frac{1}{5}\]However, by AM-GM:
 \[\cos^2\alpha+\cos^2\beta \ge 2\cos\alpha\beta\]Therefore,
 \[1-\cos^2\alpha-\cos^2\beta+\cos^2\alpha\cos^2\beta \le 1-2\cos\alpha\beta+\cos^2\alpha\cos^2\beta = (1-\cos\alpha\cos\beta)^2\]
 \[\sqrt{1-\cos^2\alpha-\cos^2\beta+\cos^2\alpha\cos^2\beta} \le 1-\cos\alpha\cos\beta\]So,
 \[\frac{1}{5} \ge \cos\alpha\cos\beta - (1-\cos\alpha\cos\beta) = 2\cos\alpha\cos\beta-1\]
 \[\frac{3}{5} \ge \cos\alpha\cos\beta\].However, the area of the rectangle is just $AS \cdot AQ = 31\cos\alpha \cdot 40\cos\beta \le 31 \cdot 40 \cdot \frac{3}{5} = \boxed{744}$.

 


 \textbf{Note on Problem Validity} It has been noted that this answer won't actually work. Let angle $QAB = m$ and angle $CAS = n$ as in Solution 1. Since we know (through that solution) that $m = n$, we can call them each $\theta$. The height of the rectangle is $AS = 31\cos\theta$, and the distance $BQ = 40\sin\theta$. We know that, if the triangle is to be inscribed in a rectangle, $AS \geq BQ$.

 
 \[AS \geq BQ\]

 \[31\cos\theta \geq 40\sin\theta\]

 \[\frac{31}{40} \geq \tan\theta\]

 However, $\tan\theta = \tan(\frac{90-A}{2}) = \frac{\sin(90-A)}{\cos(90-A)+1} = \frac{\cos A}{\sin A + 1} = \frac{\frac{2\sqrt6}{5}}{\frac{6}{5}} = \frac{\sqrt6}{3} > \frac{31}{40}$, so the triangle does not actually fit in the rectangle: specifically, B is above R and thus in the line containing segment QR but not on the actual segment or in the rectangle.

 
                    \begin{figure}[H]
                        \centering
                        \includegraphics[width=0.5\textwidth]{Solution/AIME_I_2016_9_Solution_Figure_2.png}
                    \end{figure}
                    

 The actual answer is a radical near $728$ (letting the triangle be inside the rectangle). The CAMC, however, has decided to accept only the answer $744$ despite the invalid problem statement.

 


	\item (\textit{AIME I 2016}) A strictly increasing sequence of positive integers $a_1$, $a_2$, $a_3$, $\cdots$ has the property that for every positive integer $k$, the subsequence $a_{2k-1}$, $a_{2k}$, $a_{2k+1}$ is geometric and the subsequence $a_{2k}$, $a_{2k+1}$, $a_{2k+2}$ is arithmetic. Suppose that $a_{13} = 2016$. Find $a_1$.


 
 \textbf{Answer: }504



\textbf{Solution 1}

We first create a similar sequence where $a_1=1$ and $a_2=2$. Continuing the sequence,

 
 \[1, 2,4,6,9,12,16,20,25,30,36,42,49,\cdots\]

 Here we can see a pattern; every second term (starting from the first) is a square, and every second term (starting from the third) is the end of a geometric sequence. This can be proven by induction. Similarly, $a_{13}$ would also need to be the end of a geometric sequence (divisible by a square). We see that $2016$ is $2^5 \cdot 3^2 \cdot 7$, so the squares that would fit in $2016$ are $1^2=1$, $2^2=4$, $3^2=9$, $2^4=16$, $2^2 \cdot 3^2 = 36$, and $2^4 \cdot 3^2 = 144$. By simple inspection $144$ is the only plausible square, since the other squares in the sequence don't have enough elements before them to go all the way back to $a_1$ while still staying as positive integers. $a_{13}=2016=14\cdot 144$, so $a_1=14\cdot 36=\fbox{504}$.

 



\textbf{Solution 2}

Setting $a_1 = a$ and $a_2 = ka$, the sequence becomes:

 
 \[a, ka, k^2a, k(2k-1)a, (2k-1)^2a, (2k-1)(3k-2)a, (3k-2)^2a, \cdots\]and so forth, with $a_{2n+1} = (nk-(n-1))^2a$. Then, $a_{13} = (6k-5)^2a = 2016$. Keep in mind, $k$ need not be an integer, only $k^2a, (k+1)^2a,$ etc. does. $2016 = 2^5*3^2*7$, so only the squares $1, 4, 9, 16, 36,$ and $144$ are plausible for $(6k-5)^2$. But when that is anything other than $2$, $k^2a$ is not an integer. Therefore, $a = 2016/2^2 = 504$.

 Thanks for reading, Rowechen Zhong.

 



\textbf{Solution 3}

Instead of setting $a_1$ equal to something and $a_2$ equal to something, note that it is rather easier to set $a_1=x^2$ and $a_3=y^2$ so that $a_2=xy,a_4=y(2y-x),a_5=(2y-x)^2$ and so on until you reached $a_{13}=(6y-5x)^2$ (Or simply notice the pattern), so $6y-5x=\sqrt{2016}=12\sqrt{14}$. Note that since each of the terms has degree 2 so if you multiply $x$ and $y$ by $\sqrt{14}$ you multiply each term by $14$ so each term is still a integer if the terms are already integers before you multiply $x$ and $y$ by $\sqrt{14}$, so let $w=\frac{x}{\sqrt{14}}$ and $z=\frac{y}{\sqrt{14}}$ so $6z-5w=12$. Then, for the sequence to be strictly increasing positive integers we have $(w,z)=(6,7)$ so $x=6\sqrt{14}$ and $a_1=x^2=6^2 \cdot 14=\boxed{504}$~\href{https://artofproblemsolving.com/wiki/index.php?title=Ddk001}{Ddk001}

 



\textbf{Solution 4(very risky and very stupid)}

The thirteenth term of the sequence is $2016$, which makes that fourteenth term of the sequence $2016+r$ and the $15^{\text{th}}$ term $\frac{(2016+r)^2}{2016}$. We note that $r$ is an integer so that means $\frac{r^2}{2016}$ is an integer. Thus, we assume the smallest value of $r$, which is $168$. We bash all the way back to the first term and get our answer of $\boxed{504}$.

 -Pleaseletmewin

 



\textbf{Solution 5}

Let $a_{2k-1}=s$ where $k=1$. Then, $a_{2k}=sr \Longrightarrow a_{2k+1}=sr^2$. Continuing on, we get $a_{2k+2}=sr(r-1)+sr^2=sr(2r-1) \Longrightarrow a_{2k+3}=sr^2(\frac{2r-1}{r})^2=s(2r-1)^2$. Moreover, $a_{2k+4}=s(2r-1)(r-1)+s(2r-1)^2=s(2r-1)(3r-2) \Longrightarrow a_{2k+5}=s(2r-1)^2(\frac{3r-2}{2r-1})^2=s(3r-2)^2$.

 It is clear now that $a_{2k+2c}=s(cr-(c-1))((c+1)r-c)$ and $a_{2k+2c-1}=s(cr-(c-1))^2$. Plugging in $c=6$, $a_{13}=s(6r-5)^2=2016$. The prime factorization of $2016=2^5\cdot3^2\cdot7$ so we look for perfect squares.

 $6r-5\equiv (6r-5)^2\equiv 1\pmod{6}$ if $r$ is an integer, and $\frac{\omega+5}{6}=r \Longrightarrow 6\mid{s}$ if $r$ is not an integer and $\omega$ is rational, so $6\mid{s}$. This forces $s=2\cdot3^2\cdot7\cdot{N}$. Assuming $(6r-5)$ is an integer, it can only be $2^x$, $x\in{1,2}$.

 If $6r-5=2^1$, $r=\frac{7}{6}$. If $6r-5=2^2$, $r=\frac{3}{2}$. Note that the latter cannot work since $a_{2k+1}=s(\frac{9}{4}) \Longrightarrow 4\mid{s}$ but $N=1 \Longrightarrow s=2\cdot3^2\cdot7$ in this scenario. Therefore, $r=\frac{7}{6} \Longrightarrow s=\frac{2016}{2^2}=504$. Plugging back $k=1$, $a_1=s=\boxed{504}$




\textbf{Solution 6 (Headache)}

Bash all the way from $a_1$ and $a_2$ to $a_{13}=2016$ to get to the equation $a_1\left(\frac{6a_2}{a_1}-5\right)^2=2016$ and then $a_2=2\sqrt{14a_1}+\frac{5}{6}a_1$ (direct consequence). Notice that from the two equations we can deduce that:

 1) $a_1\mid2016$,

 2) $6\mid a_1$, and

 3) $\sqrt{\frac{a_1}{14}}\in\mathbb{Z}$.

 Since $2016=2^5\cdot3^2\cdot7=14(2^4\cdot 3^2)$, the only possible values of $a_1$ that satisfy the conditions are $126$, $504$, and $2016$ itself. Obviously $2016$ cannot be a correct answer choice (from AIME rules), and if $a_1=126$, $a_3\not\in\mathbb{Z}$, so $\boxed{504}$ is the solution.

 



\textbf{Solution 7}

Let $a_{11} = x$. Then, $a_{12} = xr$ and $a_{13} = xr^2 = 2016$. Note, in order to alternate between arithmetic and geometric sequences, you would want $r$ to be as small as possible so our sequence is always positive as stated in the question. Since $2016 = 2^5 \cdot 3^2 \cdot 7$, the smallest ratio, r, (which must be greater than 1 as the sequence is strictly increasing) must satisfy that $a_12 = \frac{2016}{r}$ and $a_11 = \frac{2016}{r^2}$ are both integers. Therefore, $\frac{1}{r}$ can at min be $\frac{m}{2^2 \cdot 3}$. To minimize this we have $m = 11$ so $r = \frac{12}{11}$. Doing this and working backwards, we get the sequence$a_1 = 504 a_2 = 588 a_3 = 686 a_4 = 784 a_5 = 896 a_6 = 1008 a_7 = 1134 a_8 = 1260 a_9 = 1400 a_{10} = 1540 a_{11} = 1694 a_{12} = 1848 a_{13} = 2016$so $a_1 = \boxed{504}$~Sid-darth-vater

 



\textbf{Solution 8 (Induction)}

Let $a_1=a,a_2=b$, we know that for every positive integer $k$,\begin{align*}a_{2k}^{2}&=a_{2k-1}\cdot a_{2k+1}\\2a_{2k+1}&=a_{2k}+a_{2k+2}\end{align*}
Computing some following terms, we have\begin{align*}a_3&=\frac{b^2}{a}\\a_4&=\frac{b(2b-a)}{a}\\a_5&=\frac{(2b-a)^2}{a}\\a_6&=\frac{(2b-a)(3b-2a)}{a}\end{align*}
We see a pattern for the 2 sequences consisting of the odd terms and even terms of $\{a_{n}\}$. We proceed with induction.

 Let $c_k=a+(k-1)d$, where $d=b-a$. We attempt to prove the following:
 \[a_{2k-1}=\frac{c_{k}^{2}}{a}, \qquad a_{2k}=\frac{c_{k}c_{k+1}}{a}\]This is obviously true for $k=1$. Assuming this statement is true for a positive integer $k$, then
 \[a_{2k+1}=\frac{a_{2k}^{2}}{a_{2k-1}}=\frac{(\frac{c_{k}c_{k+1}}{a})^2}{\frac{c_{k}^2}{a}}=\frac{c_{k+1}^2}{a}\]Also, using the rule of the arithmetic subsequence, we have
 \[a_{2k+2}=2a_{2k+1}-a_{2k}=\frac{c_{k+1}(2c_{k+1}-c_{k})}{a}\]We see that
 \[2c_{k+1}-c_k=2a+2kd-a-(k-1)d=a+(k+1)d=c_{k+2}\]Thus, $a_{2k+2}=\frac{c_{k+1}c_{k+2}}{a}$, our induction hypothesis is hence proven.

 Note that $a \mid c_{k}^2 \iff a \mid d^2$. Since $a_{13}=2016$, we have $\frac{(a+6d)^2}{a}=2016$, let $k=a+6d$. This follows that
 \[a=\frac{k^2}{2016}\]Note that $2016=2^5\cdot 3^2\cdot 7$. Substituting $k=168m$, this follows that $a=14m^2$. Therefore
 \[6d=k-a=168m-14m^2\implies d=\frac{7}{3}m(12-m)\]$m$ needs to be a multiple of $3$, substituting $m=3j$, we find
 \[d=7j(12-3j)\]The cases for $j=1,2,3$ gives the solutions
 \[(a, d)=(126,63),(504,84),(1134,63)\]Note that only the second ordered-pair suffices $a \mid d^2$, the solution is hence $a_{1}=a=\boxed{504}$, as desired.

 


	\item (\textit{AIME I 2016}) Let $P(x)$ be a nonzero polynomial such that $(x-1)P(x+1)=(x+2)P(x)$ for every real $x$, and $\left(P(2)\right)^2 = P(3)$. Then $P(\frac72)=\frac{m}{n}$, where $m$ and $n$ are relatively prime positive integers. Find $m + n$.


 
 \textbf{Answer: }109



\textbf{Solution 1}

Plug in $x=1$ to get $(1-1)P(1+1) = 0 = (1+2)P(1) \Rightarrow P(1) = 0$. Plug in $x=0$ to get $(0-1)P(0+1) = (0+2)P(0)\Rightarrow P(0) = -\frac{1}{2}P(1) = 0$. Plug in $x=-1$ to get $(-1-1)P(-1+1) = (-1+2)P(-1)\Rightarrow (-2)P(0)=P(-1)\Rightarrow P(-1) = 0$. So $P(x) = x(x-1)(x+1)Q(x)$ for some polynomial $Q(x)$. Using the initial equation, once again, 
 \[(x-1)P(x+1) = (x+2)P(x)\]
 \[(x-1)((x+1)(x+1-1)(x+1+1)Q(x+1)) = (x+2)((x)(x-1)(x+1)Q(x))\]
 \[(x-1)(x+1)(x)(x+2)Q(x+1) = (x+2)(x)(x-1)(x+1)Q(x)\]
 \[Q(x+1) = Q(x)\]From here, we know that 

 


	\item (\textit{AIME I 2016}) Find the least positive integer $m$ such that $m^2 - m + 11$ is a product of at least four not necessarily distinct primes. 


 
 \textbf{Answer: }132



\textbf{Solution 1}

$m(m-1)$ is the product of two consecutive integers, so it is always even. Thus $m(m-1)+11$ is odd and never divisible by $2$. Thus any prime $p$ that divides $m^2-m+11$ must divide $4m^2-4m+44=(2m-1)^2+43$. We see that $(2m-1)^2\equiv -43\pmod{p}$. We can verify that $-43$ is not a perfect square mod $p$ for each of $p=3,5,7$. Therefore, all prime factors of $m^2-m+11$ are $\ge 11$.

 Let $m^2 - m + 11 = pqrs$ for primes $11\le p \le q \le r \le s$. From here, we could go a few different ways:

 



\textbf{Solution 1a}

Suppose $p=11$; then $m^2-m+11=11qrs$. Reducing modulo 11, we get $m\equiv  1,0 \pmod{11}$ so $k(11k\pm 1)+1 = qrs$. 

 Suppose $q=11$. Then we must have $11k^2\pm k + 1 = 11rs$, which leads to $k\equiv \pm 1 \pmod{11}$, i.e., $k\in \{1,10,12,21,23,\ldots\}$. 

 $k=1$ leads to $rs=1$ (impossible)! Then $k=10$ leads to $rs=101$, a prime (impossible). Finally, for $k=12$ we get $rs=143=11\cdot 13$. 

 Thus our answer is $m=11k= \boxed{132}$.

 



\textbf{Solution 1b}

Let $m^2 - m + 11 = pqrs$ for primes $p, q, r, s\ge11$. If $p, q, r, s = 11$, then $m^2-m+11=11^4$. We can multiply this by $4$ and complete the square to find $(2m-1)^2=4\cdot 11^4-43$. But 
 \[(2\cdot 11^2-1)^2=4\cdot 11^4-4\cdot 11^2+1 <4\cdot 11^4-43<(2\cdot 11^2)^2,\]hence we have pinned a perfect square $(2m-1)^2=4\cdot 11^4-43$ strictly between two consecutive perfect squares, a contradiction. Hence $pqrs \ge 11^3 \cdot 13$. Thus $m^2-m+11\ge 11^3\cdot 13$, or $(m-132)(m+131)\ge0$. From the inequality, we see that $m \ge 132$. $132^2 - 132 + 11 = 11^3 \cdot 13$, so $m = \boxed{132}$ and we are done.

 



\textbf{Solution 2}

First, we can show that $2,3,5,7$ $\not |$ $m^2 - m + 11$. This can be done by just testing all residue classes. 

 For example, we can test $m \equiv 0 \mod 2$  or  $m \equiv 1 \mod 2$  to show that $m^2 - m + 11$ is not divisible by 2.

 Case 1: m = 2k  

 \begin{verbatim}
	$m^2 - m + 11  \equiv 2(2 \cdot k^2 - k + 5) +1 \equiv 1 \mod 2$
\end{verbatim}

Case 2: m = 2k+1  

 \begin{verbatim}
	$m^2 - m + 11  \equiv 2(2 \cdot k^2 + k + 5) +1 \equiv 1 \mod 2$
\end{verbatim}

Now, we can test $m^2 - m + 11 = 11^4$, which fails, so we test $m^2 - m + 11 = 11^3 \cdot 13$, and we get m = $132$.

 -AlexLikeMath

 


	\item (\textit{AIME I 2016}) Freddy the frog is jumping around the coordinate plane searching for a river, which lies on the horizontal line $y = 24$. A fence is located at the horizontal line $y = 0$. On each jump Freddy randomly chooses a direction parallel to one of the coordinate axes and moves one unit in that direction. When he is at a point where $y=0$, with equal likelihoods he chooses one of three directions where he either jumps parallel to the fence or jumps away from the fence, but he never chooses the direction that would have him cross over the fence to where $y < 0$. Freddy starts his search at the point $(0, 21)$ and will stop once he reaches a point on the river. Find the expected number of jumps it will take Freddy to reach the river.


 
 \textbf{Answer: }273



\textbf{Solution 1 (Substitution)}

Clearly Freddy's $x$-coordinate is irrelevant, so we let $E(y)$ be the expected value of the number of jumps it will take him to reach the river from a given $y$-coordinate. Observe that $E(24)=0$, and 
 \[E(y)=1+\frac{E(y+1)+E(y-1)+2E(y)}{4}\] for all $y$ such that $1\le y\le 23$. Also note that $E(0)=1+\frac{2E(0)+E(1)}{3}$. This gives $E(0)=E(1)+3$. Plugging this into the equation for $E(1)$ gives that 
 \[E(1)=1+\frac{E(2)+3E(1)+3}{4},\] or $E(1)=E(2)+7$. Iteratively plugging this in gives that $E(n)=E(n+1)+4n+3$. Thus $E(23)=E(24)+95$, $E(22)=E(23)+91=E(24)+186$, and $E(21)=E(22)+87=E(24)+273=\boxed{273}$.

 



\textbf{Solution 2 (Cancellation)}

As in Solution $1$, let $E(y)$ be the expected value of jumps it takes Freddy to reach the river from the specific $y$ coordinate. We can set up a system of equations. The equation for $E(0)$ is $E(0)=\frac{2}{3}(E(0)+1)+\frac{1}{3}(E(1)+1)$, which, rearranging, gives $E(0)=E(1)+3.$ The equation for $E(23)$ is $E(23)=\frac{1}{2}(E(23)+1)+\frac{1}{4}(E(22)+1)+\frac{1}{4}$ (since $E(24)=0$), and rearranging gives $2E(23)=E(22)+4$.

 
For any other $y$ such that $0<y<23$, we can get the equation $E(y)=\frac{1}{2}(E(y)+1)+\frac{1}{4}(E(y-1)+1)+\frac{1}{4}(E(y+1)+1)$, which can be rearranged to $2E(y)=E(y-1)+E(y+1)+4$. Thus, we can get the following system of equations:\begin{align*}2E(23)&=E(22)+4\\2E(22)&=E(21)+E(23)+4\\2E(21)&=E(20)+E(22)+4\\.\\.\\.\\2E(2)&=E(1)+E(3)+4\\2E(1)&=E(0)+E(2)+4\\E(0)&=E(1)+3.\end{align*}


 Adding all the equations and cancelling results in $E(23)=3+4(23)=95$. From the first equation, $E(22)=2\cdot 95-4=186$, and from the second equation, $E(21)=2\cdot 186-95-4=\boxed{273}.$



	\item (\textit{AIME I 2016}) Centered at each lattice point in the coordinate plane are a circle radius $\frac{1}{10}$ and a square with sides of length $\frac{1}{5}$ whose sides are parallel to the coordinate axes. The line segment from $(0,0)$ to $(1001, 429)$ intersects $m$ of the squares and $n$ of the circles. Find $m + n$.


 
 \textbf{Answer: }574



\textbf{Solution 1}

First note that $1001 = 143 \cdot 7$ and $429 = 143 \cdot 3$ so every point of the form $(7k, 3k)$ is on the line.   Then consider the line $l$ from $(7k, 3k)$ to $(7(k + 1), 3(k + 1))$.  Translate the line $l$ so that $(7k, 3k)$ is now the origin.   There is one square and one circle that intersect the line around $(0,0)$.  Then the points on $l$ with an integral $x$-coordinate are, since $l$ has the equation $y = \frac{3x}{7}$:

 
 \[(0,0), \left(1, \frac{3}{7}\right), \left(2, \frac{6}{7}\right), \left(3, 1 + \frac{2}{7}\right), \left(4, 1 + \frac{5}{7}\right), \left(5, 2 + \frac{1}{7}\right), \left(6, 2 + \frac{4}{7}\right), (7,3).\]

 We claim that the lower right vertex of the square centered at $(2,1)$ lies on $l$.  Since the square has side length $\frac{1}{5}$, the lower right vertex of this square has coordinates $\left(2 + \frac{1}{10},  1 - \frac{1}{10}\right) = \left(\frac{21}{10}, \frac{9}{10}\right)$.  Because $\frac{9}{10} = \frac{3}{7} \cdot \frac{21}{10}$, $\left(\frac{21}{10}, \frac{9}{10}\right)$ lies on $l$.  Since the circle centered at $(2,1)$ is contained inside the square, this circle does not intersect $l$.  Similarly the upper left vertex of the square centered at $(5,2)$ is on $l$.  Since every other point listed above is farther away from a lattice point (excluding (0,0) and (7,3)) and there are two squares with centers strictly between $(0,0)$ and $(7,3)$ that intersect $l$.  Since there are $\frac{1001}{7} = \frac{429}{3} = 143$ segments from $(7k, 3k)$ to $(7(k + 1), 3(k + 1))$, the above count is yields $143 \cdot 2 = 286$ squares.  Since every lattice point on $l$ is of the form $(3k, 7k)$ where $0 \le k \le 143$, there are $144$ lattice points on $l$. Centered at each lattice point, there is one square and one circle, hence this counts $288$ squares and circles.  Thus $m + n = 286 + 288 = \boxed{574}$.

 (Solution by gundraja)

 



\textbf{Solution 2}

See if you can solve the problem with the following.

 
                    \begin{figure}[H]
                        \centering
                        \includegraphics[width=0.5\textwidth]{Solution/AIME_I_2016_14_Solution_Figure_1.png}
                    \end{figure}
                    




\textbf{Solution to Solution 2}

This is mostly a clarification to Solution 1, but let's take the diagram for the origin to $(7,3)$. We have the origin circle and square intersected, then two squares, then the circle and square at $(7,3)$. If we take the circle and square at the origin out of the diagram, we will be able to repeat the resulting segment (with its circles and squares) end to end from $(0,0)$ to $(1001,429)$, which forms the line we need without overlapping. Since $143$ of these segments are needed to do this, and $3$ squares and $1$ circle are intersected with each, there are $143 \cdot (3+1) = 572$ squares and circles intersected. Adding the circle and square that are intersected at the origin back into the picture, we get that there are $572+2=\boxed{574}$ squares and circles intersected in total.

 



\textbf{Solution to Solution 2 without a diagram}

This solution is a more systematic approach for finding when the line intersects the squares and circles. Because $1001 = 7*11*13$ and $429=3*11*13$, the slope of our line is $\frac{3}{7}$, and we only need to consider the line in the rectangle from the origin to $(7,3)$, and we can iterate the line $11*13=143$ times. First, we consider how to figure out if the line intersects a square. Given a lattice point $(x_1, y_1)$, we can think of representing a square centered at that lattice point as all points equal to $(x_1 \pm a, y_1 \pm b)$ s.t. $0 \leq a,b \leq \frac{1}{10}$. If the line $y = \frac{3}{7}x$ intersects the square, then we must have $\frac{y_1 + b}{x_1 + a} = \frac{3}{7}$. The line with the least slope that intersects the square intersects at the bottom right corner and the line with the greatest slope that intersects the square intersects at the top left corner; thus we must have that $\frac{3}{7}$ lies in between these slopes, or that $\frac{y_1-\frac{1}{10}}{x_1+\frac{1}{10}} \leq \frac{3}{7} \leq \frac{y_1+\frac{1}{10}}{x_1-\frac{1}{10}}$. Simplifying, $3x_1 - 1 \leq 7y_1 \leq 3x_1 + 1$. Because $y$ can only equal $0, 1, 2, 3$, we just do casework based on the values of $y$ and find that the points $(2, 1)$ and $(5, 2)$ are intersected just at the corner of the square and $(0, 0), (7, 3)$ are intersected through the center of the square. However, we disregard one of $(0, 0)$ and $(7, 3)$, WLOG $(0, 0)$, since we just use it in our count for the next of the 143 segments. Therefore, in one of our "segments", 3 squares are intersected and 1 circle is intersected giving 4 total. Thus our answer is $143*4 = 572$. HOWEVER, we cannot forget that we ignored $(0, 0)$, which contributes another square and circle to our count, making the final answer $572 + 2 = \boxed{574}$.

 -Patrick4President

 


	\item (\textit{AIME I 2016}) Circles $\omega_1$ and $\omega_2$ intersect at points $X$ and $Y$. Line $\ell$ is tangent to $\omega_1$ and $\omega_2$ at $A$ and $B$, respectively, with line $AB$ closer to point $X$ than to $Y$. Circle $\omega$ passes through $A$ and $B$ intersecting $\omega_1$ again at $D \neq A$ and intersecting $\omega_2$ again at $C \neq B$. The three points $C$, $Y$, $D$ are collinear, $XC = 67$, $XY = 47$, and $XD = 37$. Find $AB^2$.


 
 \textbf{Answer: }270



\textbf{Solution}

Using the radical axis theorem, the lines $\overline{AD}, \overline{BC}, \overline{XY}$ are all concurrent at one point, call it $F$. Now recall by Miquel's theorem in $\triangle FDC$ the fact that quadrilaterals $DAXY$ and $CBXY$ are cyclic implies $FAXB$ is cyclic as well. Denote $\omega_{3}\equiv(FAXB)$ and $Z\equiv\ell\cap\overline{FXY}$.

 Since point $Z$ lies on the radical axis of $\omega_{1},\omega_{2}$, it has equal power with respect to both circles, thus 
 \[AZ^{2}=\text{Pow}_{\omega_{1}}(Z)=ZX\cdot ZY=\text{Pow}_{\omega_{2}}(Z)=ZB^{2}\implies AZ=ZB.\] Also, notice that 
 \[AZ\cdot ZB=\text{Pow}_{\omega_{3}}(Z)=ZX\cdot ZF\implies ZY=ZF.\] The diagonals of quadrilateral $FAYB$ bisect each other at $Z$, so we conclude that $FAYB$ is a parallelogram. Let $u:=ZX$, so that $ZY=ZF=u+47$.

 Because $FAYB$ is a parallelogram and quadrilaterals $DAXY, CBXY$ are cyclic, 
 \[\angle DFX=\angle AFX=\angle BYX=\angle BCX=\angle FCX~~\text{and}~~\angle XDF=\angle XDA=\angle XYA=\angle XFB=\angle XFC\] so we have the pair of similar triangles $\triangle DFX~\sim~\triangle FCX$. Thus 
 \[\dfrac{37}{2u+47}=\dfrac{2u+47}{67}\implies 2u+47=\sqrt{37\cdot 67}\implies u=\dfrac{1}{2}\left(\sqrt{37\cdot 67}-47\right).\] Now compute 
 \[AB^{2}=4AZ^{2}=4\cdot ZX\cdot ZY=4u(u+47)=37\cdot 67-47^{2}=\textbf{270}.\]

                    \begin{figure}[H]
                        \centering
                        \includegraphics[width=0.5\textwidth]{Solution/AIME_I_2016_15_Solution_Figure_1.png}
                    \end{figure}
                    




\textbf{Solution 1}

Let $Z = XY \cap AB$. By the radical axis theorem $AD, XY, BC$ are concurrent, say at $P$. Moreover, $\triangle DXP \sim \triangle PXC$ by simple angle chasing. Let $y = PX, x = XZ$. Then 
 \[\frac{y}{37} = \frac{67}{y} \qquad \implies \qquad y^2 = 37 \cdot 67.\] Now, $AZ^2 = \frac 14 AB^2$, and by power of a point, 
 \begin{align*}         x(y-x) &= \frac 14 AB^2, \quad \text{and} \\         x(47+x) &= \frac 14 AB^2     \end{align*}
 Solving, we get 
 \[\frac 14 AB^2 = \frac 12 (y-47)\cdot \frac 12 (y+47) \qquad \implies\]
 \[\qquad AB ^ 2 = 37\cdot 67 - 47^2 = \boxed{270}\]




\textbf{Solution 2}

By the Radical Axis Theorem $AD, XY, BC$ concur at point $E$.

 Let $AB$ and $EY$ intersect at $S$. Note that because $AXDY$ and $CYXB$ are cyclic, by Miquel's Theorem $AXBE$ is cyclic as well. Thus
 \[\angle AEX = \angle ABX = \angle XCB = \angle XYB\]and
 \[\angle XEB = \angle XAB = \angle XDA = \angle XYA.\]Thus $AY \parallel EB$ and $YB \parallel EA$, so $AEBY$ is a parallelogram. Hence $AS = SB$ and $SE = SY$. But notice that $DXE$ and $EXC$ are similar by $AA$ Similarity, so $XE^2 = XD \cdot XC = 37 \cdot 67$. But
 \[XE^2 - XY^2 = (XE + XY)(XE - XY) = EY \cdot 2XS = 2SY \cdot 2SX = 4SA^2 = AB^2.\]Hence $AB^2 = 37 \cdot 67 - 47^2 = \boxed{270}.$




\textbf{Solution 3}

First, we note that as $\triangle XDY$ and $\triangle XYC$ have bases along the same line, $\frac{[\triangle XDY]}{[\triangle XYC]}=\frac{DY}{YC}$. We can also find the ratio of their areas using the circumradius area formula. If $R_1$ is the radius of $\omega_1$ and if $R_2$ is the radius of $\omega_2$, then 
 \[\frac{[\triangle XDY]}{[\triangle XYC]}=\frac{(37\cdot 47\cdot DY)/(4R_1)}{(47\cdot 67\cdot YC)/(4R_2)}=\frac{37\cdot DY\cdot R_2}{67\cdot YC\cdot R_1}.\]Since we showed this to be $\frac{DY}{YC}$, we see that $\frac{R_2}{R_1}=\frac{67}{37}$.

 We extend $AD$ and $BC$ to meet at point $P$, and we extend $AB$ and $CD$ to meet at point $Q$ as shown below.
                    \begin{figure}[H]
                        \centering
                        \includegraphics[width=0.5\textwidth]{Solution/AIME_I_2016_15_Solution_Figure_2.png}
                    \end{figure}
                    As $ABCD$ is cyclic, we know that $\angle BCD=180-\angle DAB=\angle BAP$. But then as $AB$ is tangent to $\omega_2$ at $B$, we see that $\angle BCD=\angle ABY$. Therefore, $\angle ABY=\angle BAP$, and $BY\parallel PD$. A similar argument shows $AY\parallel PC$. These parallel lines show $\triangle PDC\sim\triangle ADY\sim\triangle BYC$. Also, we showed that $\frac{R_2}{R_1}=\frac{67}{37}$, so the ratio of similarity between $\triangle ADY$ and $\triangle BYC$ is $\frac{37}{67}$, or rather
 \[\frac{AD}{BY}=\frac{DY}{YC}=\frac{YA}{CB}=\frac{37}{67}.\]We can now use the parallel lines to find more similar triangles. As $\triangle AQD\sim \triangle BQY$, we know that
 \[\frac{QA}{QB}=\frac{QD}{QY}=\frac{AD}{BY}=\frac{37}{67}.\]Setting $QA=37x$, we see that $QB=67x$, hence $AB=30x$, and the problem simplifies to finding $30^2x^2$. Setting $QD=37^2y$, we also see that $QY=37\cdot 67y$, hence $DY=37\cdot 30y$. Also, as $\triangle AQY\sim \triangle BQC$, we find that
 \[\frac{QY}{QC}=\frac{YA}{CB}=\frac{37}{67}.\]As $QY=37\cdot 67y$, we see that $QC=67^2y$, hence $YC=67\cdot30y$.

 Applying Power of a Point to point $Q$ with respect to $\omega_2$, we find
 \[67^2x^2=37\cdot 67^3 y^2,\]or $x^2=37\cdot 67 y^2$. We wish to find $AB^2=30^2x^2=30^2\cdot 37\cdot 67y^2$.

 Applying Stewart's Theorem to $\triangle XDC$, we find
 \[37^2\cdot (67\cdot 30y)+67^2\cdot(37\cdot 30y)=(67\cdot 30y)\cdot (37\cdot 30y)\cdot (104\cdot 30y)+47^2\cdot (104\cdot 30y).\]We can cancel $30\cdot 104\cdot y$ from both sides, finding $37\cdot 67=30^2\cdot 67\cdot 37y^2+47^2$. Therefore, 
 \[AB^2=30^2\cdot 37\cdot 67y^2=37\cdot 67-47^2=\boxed{270}.\]




\textbf{Solution 4}


                    \begin{figure}[H]
                        \centering
                        \includegraphics[width=0.5\textwidth]{Solution/AIME_I_2016_15_Solution_Figure_3.png}
                    \end{figure}
                    First of all, since quadrilaterals $ADYX$ and $XYCB$ are cyclic, we can let $\angle DAX = \angle XYC = \theta$, and $\angle XYD = \angle CBX = 180 - \theta$, due to the properties of cyclic quadrilaterals. In addition, let $\angle BAX = x$ and $\angle ABX = y$. Thus, $\angle ADX = \angle AYX = x$ and $\angle XYB = \angle XCB = y$. Then, since quadrilateral $ABCD$ is cyclic as well, we have the following sums:
 \[\theta + x +\angle XCY + y = 180^{\circ}\]
 \[180^{\circ} - \theta + y + \angle XDY + x = 180^{\circ}\]Cancelling out $180^{\circ}$ in the second equation and isolating $\theta$ yields $\theta = y + \angle XDY + x$. Substituting $\theta$ back into the first equation, we obtain
 \[2x + 2y + \angle XCY + \angle XDY = 180^{\circ}\]Since
 \[x + y +\angle XAY + \angle XCY + \angle DAY = 180^{\circ}\]
 \[x + y + \angle XDY + \angle XCY + \angle DAY = 180^{\circ}\]we can then imply that $\angle DAY = x + y$. Similarly, $\angle YBC = x + y$. So then $\angle DXY = \angle YXC = x + y$, so since we know that $XY$ bisects $\angle DXC$, we can solve for $DY$ and $YC$ with Stewart’s Theorem. Let $DY = 37n$ and $YC = 67n$. Then
 \[37n \cdot 67n \cdot 104n + 47^2 \cdot 104n = 37^2 \cdot 67n + 67^2 \cdot 37n\]
 \[37n \cdot 67n + 47^2 = 37 \cdot 67\]
 \[n^2 = \frac{270}{2479}\]Now, since $\angle AYX = x$ and $\angle BYX = y$, $\angle AYB = x + y$. From there, let $\angle AYD = \alpha$ and $\angle BYC = \beta$. From angle chasing we can derive that $\angle YDX = \angle YAX = \beta - x$ and $\angle YCX = \angle YBX = \alpha - y$. From there, since $\angle ADX = x$, it is quite clear that $\angle ADY = \beta$, and $\angle YAB = \beta$ can be found similarly. From there, since $\angle ADY = \angle YAB = \angle BYC = \beta$ and $\angle DAY = \angle AYB = \angle YBC = x + y$, we have $AA$ similarity between $\triangle DAY$, $\triangle AYB$, and $\triangle YBC$. Therefore the length of $AY$ is the geometric mean of the lengths of $DA$ and $YB$ (from $\triangle DAY \sim \triangle AYB$). However, $\triangle DAY \sim \triangle AYB \sim \triangle YBC$ yields the proportion $\frac{AD}{DY} = \frac{YA}{AB} = \frac{BY}{YC}$; hence, the length of $AB$ is the geometric mean of the lengths of $DY$ and $YC$. We can now simply use arithmetic to calculate $AB^2$. 
 \[AB^2 = DY \cdot YC\]
 \[AB^2 = 37 \cdot 67 \cdot \frac{270}{2479}\]
 \[AB^2 = \boxed{270}\]

 \textbf{-Solution by TheBoomBox77}

 



\textbf{Solution 5 (not too different)}

Let $E = DA \cap CB$. By Radical Axes, $E$ lies on $XY$. Note that $EAXB$ is cyclic as $X$ is the Miquel point of $\triangle EDC$ in this configuration.

 Claim. $\triangle DXE \sim \triangle EXC$Proof. We angle chase. 
 \[\measuredangle XEC = \measuredangle XEB = \measuredangle XAB = \measuredangle XDA = \measuredangle XDE\]and
 \[\measuredangle XCE = \measuredangle XCB = \measuredangle XBA = \measuredangle XEA = \measuredangle XED. \square\]

 Let $F = EX \cap AB$. Note 
 \[FA^2 = FX \cdot FY = FB^2\]and
 \[EF \cdot FX = AF \cdot FB = FA^2 = FX \cdot FY \implies EF = FY\]By our claim, 
 \[\frac{DX}{XE} = \frac{EX}{XC} \implies EX^2 = DX \cdot XC = 67 \cdot 37 \implies FY = \frac{EY}{2} = \frac{EX+XY}{2} = \frac{\sqrt{67 \cdot 37}+47}{2}\]and
 \[FX=FY-47=\frac{\sqrt{67 \cdot 37}-47}{2}\]Finally, 
 \[AB^2 = (2 \cdot FA)^2 = 4 \cdot FX \cdot FY = 4 \cdot \frac{(67 \cdot 37) - 47^2}{4} = \boxed{270}. \blacksquare\]~Mathscienceclass

 


 



\textbf{Solution 6 (No words)}


                    \begin{figure}[H]
                        \centering
                        \includegraphics[width=0.5\textwidth]{Solution/AIME_I_2016_15_Solution_Figure_4.png}
                    \end{figure}
                    
                    \begin{figure}[H]
                        \centering
                        \includegraphics[width=0.5\textwidth]{Solution/AIME_I_2016_15_Solution_Figure_5.png}
                    \end{figure}
                    $AB^2 = 4 AM^2 =2x(2x+ 2 XY) =(XP - XY) (XP + XY) = XP^2 - XY^2 = XC \cdot XD - XY^2 = 67 \cdot 37 - 47^2 = \boxed{270}.$

 \textbf{vladimir.shelomovskii@gmail.com, vvsss}

 



\textbf{Solution 7 (Linearity of Power of a Point)}

Extend $\overline{AD}$ and $\overline{BC}$ to meet at point $P$. Let $M$ be the midpoint of segment $AB$. Then by radical axis on $(ADY)$, $(BCY)$ and $(ABCD)$, $P$ lies on $XY$. By the bisector lemma, $M$ lies on $XY$. It is well-known that $P$, $A$, $X$, and $B$ are concyclic. By Power of a point on $M$ with respect to $(PAXB)$ and $(ADY)$, 
 \[|\text{Pow}(M, (PAXB))| = MX \cdot MP = MA^2 = |\text{Pow}(M, (ADY))| = MX \cdot MY,\] so $MP=MY$. Thus $AB$ and $PY$ bisect each other, so $PAYB$ is a parallelogram. This implies that 
 \[\angle DAY = \angle YBC,\] so by the inscribed angle theorem $\overline{XY}$ bisects $\angle DXC$. 

 Claim: $AB^2 = DY \cdot YC$.

 Proof. Define the linear function $f(\bullet) := \text{Pow}(\bullet, (ADY)) - \text{Pow}(\bullet, (ABCD))$. Since $\overline{BY}$ is parallel to the radical axis $\overline{AD}$ of $(ADY)$ and $(ABCD)$ by our previous parallelism, $f(B)=f(Y)$. Note that $f(B)=AB^2$ while $f(Y)=DY \cdot YC$, so we  conclude. $\square$

 By Stewart's theorem on $\triangle DXC$, $DY \cdot YC=37 \cdot 67 - 47^2 = 270$, so $AB^2=\boxed{270}$. 

 


\end{enumerate}